% Generated draft: requires mathematical and source review.
\documentclass[11pt]{article}
\usepackage[margin=25mm]{geometry}
\usepackage{fontspec}
\setmainfont{Latin Modern Roman}
\usepackage{amsmath,amssymb,mathtools,unicode-math}
\setmathfont{Latin Modern Math}
\usepackage{enumitem,needspace,xurl,hyperref}
\setlength{\parindent}{0pt}
\setlength{\parskip}{5pt}
\setlength{\emergencystretch}{3em}
\newcommand{\N}{\mathbb{N}}
\newcommand{\Z}{\mathbb{Z}}
\newcommand{\Q}{\mathbb{Q}}
\newcommand{\R}{\mathbb{R}}
\newcommand{\C}{\mathbb{C}}
\newcommand{\F}{\mathbb{F}}
\newcommand{\A}{\mathbb{A}}
\newcommand{\PP}{\mathbb{P}}
\newcommand{\E}{\mathbb{E}}
\newcommand{\eps}{\varepsilon}
\newcommand{\dd}{\,\mathrm{d}}
\begin{document}
\section*{UnsolvedMath ID: 30004564}
\textbf{Source problem code:} catalogue entry
\subsection{Definitions and mathematical statement}
Let \((M,g)\) denote the Riemannian manifold in the supplied question. Write \(s_{\mathrm{ret}}\) and \(s_{\mathrm{int}}\) symbolically for, respectively, a retracted thin plate spline and an intrinsic manifold thin plate spline constructed from the same data, whenever both constructions are defined. These symbols identify the objects to be compared; the record does not specify their defining variational problems or the retraction construction.

\textbf{Comparison problem.} Determine how \(s_{\mathrm{ret}}\) differs from \(s_{\mathrm{int}}\), particularly in the interpolation regime denoted in the source by \(\lambda=\infty\), and describe how this comparison depends on the manifold's closeness to Euclidean space. The supplied background also identifies comparison of smoothness properties as an unresolved target.

\textbf{Unresolved specifications.} The record does not specify the parameter domain, interpolation sites and values, admissible maps, spline energies, boundary conditions, choice of retraction, existence or uniqueness assumptions, or a measure of difference between the splines. It also does not define closeness to Euclidean space or the smoothness properties to be compared. The notation \(\lambda=\infty\) is supplied as a designation of the interpolation case; no penalized functional or limiting procedure is given. Consequently, no particular quantitative estimate, convergence statement, or universal quantification over manifolds and data can be extracted without additional information.

The supplied dated assessment says that the interpolation difference has been quantified under \(\varepsilon\)-flatness, while a general smoothness comparison remains unsettled. Neither \(\varepsilon\)-flatness nor the hypotheses and conclusion of that reported comparison are supplied, so this draft does not assert such a result.
\subsection{Short English statement}
Compare retracted and intrinsic thin plate splines on a Riemannian manifold, especially for interpolation, denoted by \(\lambda=\infty\), and determine how their difference and smoothness depend on closeness to Euclidean geometry. Precise spline constructions, comparison criteria, and geometric hypotheses remain unspecified in the supplied record.
\subsection{Sources}
\begin{itemize}
\item Local UnsolvedMath registry, \texttt{data/unsolvedmath/problems.json}, permanent ID \(30004564\), source code \texttt{catalogue entry}, title: Retracted versus Intrinsic Thin Plate Splines on Manifolds. The supplied record belongs to the Oberwolfach Reports Open Problems collection. No individual proposer is identified; the proposed year is recorded as 2020.
\item Supplied source citation: \emph{Mini-Workshop: Computational Optimization on Manifolds (2021)}, Oberwolfach Reports, pp.\ 1791--1821, DOI 10.4171/owr/2020/36. \url{https://doi.org/10.4171/owr/2020/36}. The background describes Oberwolfach Reports as records of mathematical workshops and their open-problem discussions.
\item Supplied literature source: Steinke et al., manifold-valued thin-plate splines. \url{https://doi.org/10.1111/j.1467-8659.2008.01141.x}. No fuller bibliographic details or theorem statement are supplied.
\item Supplied literature source: OWR 2020/36 report. \url{https://doi.org/10.4171/owr/2020/36}.
\item Literature assessment dated exactly \texttt{2026-08-22}: supplied dated curation, not independently verified. The record labels the problem partially solved and reports that manifold thin-plate splines and retraction-based splines are established, that their interpolation difference is quantified under \(\varepsilon\)-flatness, and that a general comparison of smoothness properties remains unsettled. The background explicitly describes this as dated literature triage derived from the supplied OWR curation and requests independent verification of time-sensitive status and cited claims.
\item The registry supplies the statement-status label \texttt{corrected\_verified} and reports a manual semantic, typographic, and LaTeX review claiming that the clean statement is self-contained and preserves the mathematical target. These are supplied registry qualifications, not independent verification by this draft.
\end{itemize}
\end{document}
